%! Author = denis
%! Date = 11/08/2026

\begin{frame}{Códigos Digitais}
    \small

    \begin{block}{Definição}
        Em sistemas digitais, diferentes códigos podem ser utilizados
        para representar números, caracteres e informações.
    \end{block}

    \vspace{0.2cm}

    Neste conteúdo estudaremos três códigos importantes:

    \begin{itemize}
        \item \textbf{BCD} -- representação de algarismos decimais em binário;
        \item \textbf{Gray} -- representação em que apenas um bit muda entre valores consecutivos;
        \item \textbf{ASCII} -- representação numérica de caracteres.
    \end{itemize}

    \vspace{0.3cm}

    \begin{alertblock}{Importante}
        Apesar de utilizarem bits, BCD, Gray e ASCII possuem
        \textbf{objetivos diferentes}.
    \end{alertblock}

\end{frame}